\scanpage{88}
% archive foliation: 87
\deleted{\unsure{Splittages|Splittage de}}

\uline{1.3.}\quad fibré \unsure{projectif en droites} sur $X$ avec
\uline{deux} sections \uline{disjointes} $\sigma,\sigma'$.

\diagram{\begin{tikzcd}[row sep=small] {}\arrow[d,leftrightarrow]\\{}\end{tikzcd}}
fibré inversible $\Omega$ \quad
$(P=\mathbb P(\Omega\oplus\underline O_X)\simeq\overline{W(\Omega)}$
\margin{clôture projective\quad $)$}
\margin{avec $\sigma=$ section à l'infini\par $\sigma'=$ section nulle}

\diagram{\begin{tikzcd}[row sep=small] {}\arrow[d,leftrightarrow]\\{}\end{tikzcd}}
[$M$ somme directe de deux Idéaux spéciaux $\ldots$]

\diagram{\begin{tikzcd}[row sep=small] {}\arrow[d,leftrightarrow]\\{}\end{tikzcd}}
$(\mathfrak g$ avec
\inserted{$\mathcal L$ et}
\deleted{\unsure{filtration scindée|deux filtrations de $\mathfrak g$}},
\unsure{l'ensemble sur 1|l'une sur l'autre}
(\unsure{cf. Ch. Klein|cf. Ch. Killing})
\par soit décomposition en somme
\[
\mathfrak g\simeq
\underbrace{\underset{\substack{\Vert\\\mathcal L_u}}{\Omega}
+\underset{\substack{\simeq\\\underline O_X}}{\unsure{\underline O_X^{0}|\Omega^{0}}}}_{\mathcal L}
+\Omega^{-1}
\]
\[
\begin{aligned}
[\omega_1,\lambda]&=-[\lambda,\omega_1]=-\lambda\omega_1\\
[\omega_{-1},\lambda]&=-[\lambda,\omega_{-1}]=\lambda\omega_{-1}\\
[\omega_1,\omega_{-1}]&=-[\omega_{-1},\omega_1]=2\omega_1\omega_{-1}
\end{aligned}
\]
\margin{\deleted{$\unsure{\mathfrak g_{\mathcal L}|\mathfrak g_{\mathcal L_i}}$}
$\longrightarrow\unsure{\mathfrak g_{\mathcal L'}|\mathfrak g'_{\mathcal L'}}$\par
\unsure{équiv.}\par \unsure{sauf|dans} \unsure{les|en} Car. $\ne2$\unsure{!|?}\par
\uline{\unsure{essentielles|inutiles}}}
\figure{\begin{tikzpicture}[x=1cm,y=1cm]
\draw (0.3,0)--(0.05,-2.5) .. controls (-0.1,-2.8) and (-0.25,-2.9) .. (-0.4,-3)
.. controls (0,-3.1) and (-0.1,-3.3) .. (0,-4)
--(0.4,-7) .. controls (0.5,-7.4) and (0.5,-7.6) .. (0.8,-7.9);
\end{tikzpicture}}
soit $\mathfrak g$ avec $\mathfrak t\subset\mathfrak g$ sous-Alg. de
Cartan, et le choix d'une des deux sous-Algèbres \unsure{unipotentes|propres}
$\Omega,\Omega'$ \unsure{normalisées|invariantes} par $\uline{\mathfrak t}$
(ce qui fixe l'iso $\mathfrak t\simeq O_X\simeq\underline{\operatorname{Hom}}_{O_X}(\Omega,\Omega)$)

\diagram{\begin{tikzcd}[row sep=small] {}\arrow[d,leftrightarrow]\\{}\end{tikzcd}}
$\mathfrak g'$ avec données analogues
(Mais dans la table de multiplication
dans ces derniers cas il faut
\[
\begin{array}{r}
-2\lambda\omega_1,\\
2\lambda\omega_{-1}\\
\omega_1\omega_{-1}
\end{array}
\qquad)
\]
\diagram{\begin{tikzcd}[row sep=small] {}\arrow[d,leftrightarrow]\\{}\end{tikzcd}}
$G$ avec couple de \unsure{Killing}.

\diagram{\begin{tikzcd}[row sep=small] {}\arrow[d,leftrightarrow]\\{}\end{tikzcd}}
$G'$ avec couple de \unsure{Killing}.

